\section*{Avant-propos}
\addcontentsline{toc}{section}{Avant-propos}
\markboth{Avant-propos}{Avant-propos}

L'École nationale de la Statistique et de l'Analyse économique (ENSAE) forme
les cadres appelés à produire et à analyser l'information statistique
nécessaire à la conduite des politiques publiques. La filière des Ingénieurs
statisticiens économistes y prépare à concevoir des dispositifs de mesure, à
traiter des données d'enquête et à en tirer des enseignements utiles à la
décision.

Le cycle se referme sur un mémoire de fin d'études. L'exercice demande à
l'élève de choisir une question de recherche, d'en construire le cadre
d'analyse, de mobiliser des données réelles et d'en tirer des conclusions
argumentées, en assumant les choix méthodologiques qui les fondent comme les
limites qui les bornent. Il ne s'agit pas seulement de restituer des
techniques apprises, mais de les mettre à l'épreuve d'une question qui n'a pas
de réponse écrite d'avance.

Le présent mémoire s'inscrit dans cette démarche. Il porte sur l'impact des
transferts de migrants sur la pauvreté multidimensionnelle des enfants au
Sénégal, et repose sur les deux éditions de l'Enquête harmonisée sur les
conditions de vie des ménages produites par l'Agence nationale de la
Statistique et de la Démographie. Le sujet se situe au croisement de deux
préoccupations, celle de la mesure du bien-être des enfants au-delà du seul
revenu, et celle de l'évaluation d'impact sur données observationnelles, où
la difficulté ne tient pas au calcul mais à la construction d'une comparaison
qui ait un sens.

Le travail a été mené sous la direction du Docteur Mamadou Abdoulaye Diallo.
Les analyses présentées ici, ainsi que les opinions qu'elles expriment,
n'engagent que leur auteur.
